Prediction uses utilization, L2 counters, and L1/CPU counters to prediction utilization in lower CPU frequency. (e.g., if suffering from contention, can decrease frequency with low impact into performance)

Problem at ANN entry:
- L2 counters are global;
- L1/CPU counters are invariable on job if code is the same, even if the utilization changes;

therefore, a CPU Hungry task running parallel to a bandwidth will have high L2 counters. Nonetheless, dividing by CPU counters will let ANN to ignore L2 counters.

A disparity running in parallel to a bandwidth will have high L2 counters, and we like that. However, CPU counters will be the same as the ones not running in parallel to a bandwidth.

---

A stronger approach would be to use tasksets with temporal variation. Over time, the tasks would show a repeating behavior:

T0: |###       |##        |#####     |###       |##        |
T1: |######    |###       |######    |###       |######    |
T2: |##        |##        |#         |##        |##        |

The, the model would need to learn the behavior of each task. This cold be done with a global model and a local layer for each task.

Each clock level would also lead to a different behavior. This non-linear scenarior could be learned by an ANN:

(Same scenario as above, but with a lower clock):

T0: |####      |####      |######    |####      |####      |
T1: |#######   |#####     |#######   |#####     |#######   |
T2: |####      |####      |###       |####      |####      |

The ANN would take as input the clock and PMU metrics.
